\section{A polynomial unrestricted error from two girth scales}
\label{sec:two-scale-polynomial-error}

This proof combines Fable's nested-cycle argument with the exponential cut estimate.
It uses two nested short-cycle deletions.
The larger deletion finds a long chain.
Restoring the medium-length deleted cycles makes most chain blocks suitable for a spectral test.
The proof is mathematical and has not been formalized in Lean.

All logarithms in this section have base two unless written as $\ln$.
Fix
\[
 A=6e^{168},\qquad t=\lceil10000A\rceil,
 \qquad \kappa=\frac1{10^6t2^t}>0.
\]
All thresholds below are absolute.
They are not optimized.

\begin{theorem}[A polynomial forest error]
\label{thm:two-scale-polynomial-forest}
Let $\Gamma$ be a connected multigraph of rank $r$ and maximum degree at most three.
Let $P$ be a nonempty vertex set such that $\Gamma[P]$ is connected,
\[
 |P|\le\frac{2r}{\log^2r},
\]
and every vertex outside $P$ has degree three.
For a uniform even subgraph $\xi$, let $\mathcal F$ mean that $\xi$ is a forest on $\Gamma-P$.
For all sufficiently large $r$,
\[
 \Pr(\mathcal F)\le\frac12+r^{-\kappa}.
\]
\end{theorem}

\subsection{The two cycle families and their moments}

Assume, for a contradiction, that
\begin{equation}\label{eq:two-scale-assumption}
 \Pr(\mathcal F)>\frac12+\varepsilon,
 \qquad \varepsilon=r^{-\kappa}.
\end{equation}
Put
\[
 x=\log r,\qquad L=\lfloor Ax\rfloor,
 \qquad D_0=\lfloor x/100\rfloor,
 \qquad D_1=\lfloor x/10\rfloor,
 \qquad s_j=\lfloor(D_j-1)/2\rfloor.
\]
First choose a maximal disjoint family $\mathcal Q_0$ of cycles of length at most $D_0$ in $\Gamma-P$.
Extend this family to a maximal disjoint family $\mathcal Q_1$ of cycles of length at most $D_1$.
Let $S_j$ be the union of the vertices of $\mathcal Q_j$, and put
\[
 J_j=\Gamma-(P\cup S_j).
\]
Thus $S_0\subseteq S_1$ and $J_1\subseteq J_0$.

Every chosen cycle is a cyclic region with boundary at most its length.
Theorem~\ref{thm:sharper-exponential-cuts} and \eqref{eq:two-scale-assumption} give
\[
 |\mathcal Q_j|<8D_j2^{D_j}\ln(2/\varepsilon).
\]
Consequently,
\begin{equation}\label{eq:two-scale-deletion-sizes}
 |S_0|=O(x^3r^{1/100}),\qquad
 |S_1|=O(x^3r^{1/10}).
\end{equation}
Each $J_j$ is simple and has girth greater than $D_j$.
Its order $v_j$ and degree deficit $d_j$ satisfy
\[
 v_j=2r+O(|P|+|S_j|),\qquad
 d_j\le3(|P|+|S_j|),\qquad Ld_j=o(v_j).
\]

Let $\mathcal C_j$ be the cycles of $J_j$ of length at most $L$.
Write
\[
 a_C=2^{-|C|},\qquad \mu_j=\sum_{C\in\mathcal C_j}a_C.
\]
The manuscript's walk bound gives, uniformly over vertices $v$ of $J_j$,
\begin{equation}\label{eq:two-scale-local-mass}
 \sum_{C\in\mathcal C_j:v\in C}a_C\le L2^{-s_j}.
\end{equation}
The spectral supply argument, applied to $J_1$, gives
\begin{equation}\label{eq:two-scale-initial-mass}
 \mu_1\ge\left(\frac18-o(1)\right)\ln(A/6)>20.
\end{equation}
Here $L^2 2^{-s_1}=o(1)$ controls repeated non-backtracking walks, and $Ld_1=o(v_1)$ controls the spectral loss.

We record the moment bounds at both girth scales.
For disjoint $C,C'\in\mathcal C_j$, the manuscript's link decomposition now gives
\[
 \ell(C,C')\le1+\frac{2L}{t}
   +8t2^t\ln(2/\varepsilon)+(8L/D_j)^2.
\]
The first term uses the connectedness of $P$.
The third term uses the exponential cut estimate for the cyclic complementary regions.
The tree-region argument remains valid because $2t+D_j/2\le D_j$ for large $r$.
The choice of constants gives
\begin{equation}\label{eq:two-scale-link-bound}
 \ell(C,C')\le x/1000
\end{equation}
for large $r$, at both scales.
Indeed, $2L/t\le x/5000$, the third term equals
$(8\ln2/10^6)x+O_t(1)$, and the last term is bounded independently of $r$.

Use the manuscript's normalized cycle-component variables
\[
 X_C=\frac{a_C}{\Pr(C\text{ is a component of }\xi)}
       \mathbf1_{\{C\text{ is a component of }\xi\}},
 \qquad Z_j=\sum_{C\in\mathcal C_j}X_C.
\]
They satisfy $\mathbb EX_C=a_C$, $\mathbb EX_C^2\le a_C$, and
\[
 \mathbb E(X_CX_{C'})=a_Ca_{C'}2^{\ell(C,C')-1}
\]
for disjoint cycles.
Distinct intersecting cycles have zero joint contribution.
The exceptional-partner lemma and \eqref{eq:two-scale-link-bound} bound the contribution of partners with at least three links by
\[
 a_C\,2^{x/1000}L^3 2^{-s_j}=o(a_C).
\]
This holds even at the smaller scale, since $s_0=x/200+O(1)$.

Let $\Lambda_j$ be the total weight of ordered disjoint one-link pairs in $\mathcal C_j$.
There is a quantity $\eta_r\to0$ such that
\begin{equation}\label{eq:two-scale-inverse-moment}
 \mathbb EZ_j^2
 \le2\mu_j^2-\Lambda_j+(1+\eta_r)\mu_j.
\end{equation}
On $\mathcal F$, one has $Z_j=0$.
Cauchy--Schwarz and \eqref{eq:two-scale-assumption} therefore imply
\begin{equation}\label{eq:two-scale-two-consequences}
 \Lambda_j<(1+\eta_r)\mu_j,
 \qquad
 \mu_j\le\frac{1+\eta_r}{4\varepsilon}.
\end{equation}

\subsection{A chain exists after weighted averaging}

For disjoint cycles $C,C'\in\mathcal C_1$, call $C'$ \emph{far from $P$ behind $C$}
when it lies in a component of $\Gamma-V(C)$ that avoids $P$.
This relation is antisymmetric.
Indeed, a path from $P$ to its first vertex of $C$ avoids that far component.
It therefore avoids $C'$.
Since $C$ is connected and disjoint from $C'$, all of $C$ lies in the $P$-component of $\Gamma-V(C')$.
It follows that the total weight of ordered far pairs is at most $\mu_1^2/2$.

By \eqref{eq:two-scale-local-mass}, the ordered weight of intersecting pairs is at most
$\mu_1L^2 2^{-s_1}=o(\mu_1)$.
The same bound applies to disjoint pairs joined by an edge.
By \eqref{eq:two-scale-two-consequences}, one-link pairs have weight at most $(1+o(1))\mu_1$.
Thus the ordered weight of the remaining pairs is at least
\[
 \frac12\mu_1^2-(1+o(1))\mu_1.
\]
In every remaining pair, $C'$ lies in the $P$-component of $\Gamma-V(C)$,
no edge joins the two cycles, and they have at least two links.
Equation~\eqref{eq:two-scale-initial-mass} therefore supplies one cycle $C$ whose such partners have total weight at least eight.

Put $H=\Gamma-V(C)$ and let $N(C)$ be the vertices outside $C$ adjacent to $C$.
Since $C$ lies outside $P$, every vertex of $C$ is cubic, so $|N(C)|\le|C|\le L$.
For each selected partner $C'$, at least two components of $H-V(C')$ are adjacent to both cycles.
At most one contains $P$.
Choose another such component $Y$ and a vertex $u\in Y\cap N(C)$.
The component $Y$ lies in the main component of $H$, since it is adjacent to $C'$ there.
Thus $C'$ separates $u$ from $P$ in $H$.

For some fixed $u\in N(C)$, cycles separating $u$ from $P$ in $H$ have total weight at least $8/L$.
Choose a maximal vertex-disjoint subfamily of these cycles.
If it has $T$ members, its vertex union meets every cycle in the family.
Equation~\eqref{eq:two-scale-local-mass} gives
\[
 \frac8L\le TL^2 2^{-s_1},\qquad
 T\ge\frac{8\,2^{s_1}}{L^3}
 \ge\frac{r^{1/20}}{L^3}
\]
for large $r$.

These disjoint cycles are nested between $u$ and $P$ in $H$.
Here is the required separation argument.
For two such cycles, replace the middle of a $u$--$P$ path by an arc of the first cycle.
The resulting path meets the second cycle in either the $u$-component or the $P$-component after deleting the first.
Connectedness places the entire second cycle in that component.
The two cycles cannot each lie on the other's $u$-side, since a path through the first cycle would then avoid the second separator.
They cannot each lie on the other's $P$-side, by the reversed argument.
Thus they have a consistent linear order.

Write the ordered cycles as $C_1,\ldots,C_T$.
Let $U_i$ be the component of $u$ in $H-V(C_i)$.
Then $U_i\subseteq U_{i+1}$, and $C_i\subseteq U_{i+1}$.

\subsection{Restore medium cycles inside the chain}

Put $q=\lceil20000Ax\rceil+2$.
Choose disjoint index intervals $[a,b]$ with $b-a=q$, leaving a separator cycle between consecutive intervals.
There are
\[
 m_* = \Omega\!\left(\frac{r^{1/20}}{x^4}\right)
\]
such intervals.
For each interval define the full vertex set
\[
 V(a,b)=U_b\setminus\bigl(U_a\cup V(C_a)\bigr).
\]
It contains $C_{a+1},\ldots,C_{b-1}$.
Consequently,
\begin{equation}\label{eq:two-scale-block-volume}
 |V(a,b)|\ge(q-1)(D_1+1)\ge2000Ax^2.
\end{equation}
These sets are pairwise disjoint and pairwise non-adjacent in $H$.
Every edge of $H$ leaving one set ends on $C_a$ or $C_b$.
Its boundary in $H$ is therefore at most $2L$.
The set has no neighbor in $P$, since $C_b$ separates it from $P$.

The definition uses full vertex sets in $H$.
It restores the vertices in $S_1\setminus S_0$ that lie between the separator cycles.
This restoration is the reason for using two girth scales.

Discard every block meeting $N(C)$.
At most $L$ blocks are lost.
Also discard every block meeting $S_0$ or a neighbor of $S_0$ in $\Gamma$.
At most $4|S_0|$ further blocks are lost, since $\Gamma$ has maximum degree three.
By \eqref{eq:two-scale-deletion-sizes}, this loss is $o(m_*)$.
Let $V_1,\ldots,V_m$ be the remaining blocks.
Then
\begin{equation}\label{eq:two-scale-many-blocks}
 m=\Omega\!\left(\frac{r^{1/20}}{x^4}\right).
\end{equation}

Each $V_i$ lies in $J_0$.
It has no neighbor in $P\cup S_0$, so every one of its vertices has degree three in $J_0$.
Thus its degree deficit in $J_0$ is exactly zero.
Its boundary in $J_0$ has size at most $2L$.
The blocks remain pairwise non-adjacent.

\subsection{Many blocks give many spectral modes}

We include the spectral argument for completeness.
Let $J$ be a simple graph of maximum degree three.
Suppose $V_1,\ldots,V_m$ are pairwise disjoint and pairwise non-adjacent vertex sets.
Write $b_i=|\delta_J V_i|$ and
$\delta_i=\sum_{v\in V_i}(3-\deg_J v)$.
If $0<\theta<1/2$ and
\[
 \delta_i+3b_i\le\theta|V_i|,
\]
then the non-backtracking matrix of $J$ has at least $m$ real eigenvalues in $[2-\theta,2]$, with multiplicity.

To prove this, let $A_J$ be its adjacency matrix and $\Delta_J$ its degree matrix.
Put $M(\lambda)=\lambda^2I-\lambda A_J+\Delta_J-I$.
For the indicator vector of $V_i$,
\[
 \mathbf1_{V_i}^{\mathsf T}M(2-\theta)\mathbf1_{V_i}
 =-\theta(1-\theta)|V_i|
 +(1-\theta)\delta_i+(2-\theta)b_i\le0.
\]
Nonadjacency makes the corresponding cross terms zero.
Hence $M(2-\theta)$ has at least $m$ nonpositive eigenvalues.
The matrix $M(2)$ is positive semidefinite.
Its ordered eigenvalues vary continuously, so at least $m$ zeros occur in $[2-\theta,2]$, counted with kernel multiplicity.
Equivalently, $M'(\lambda)=2\lambda I-A_J$ is positive definite for $\lambda>3/2$, so each ordered eigenvalue crosses zero at most once there.
The manuscript's map
$g\mapsto f$, with $f(v\to w)=\lambda g(w)-g(v)$,
injects each kernel into the corresponding non-backtracking eigenspace.
This proves the eigenvalue count.

For even $\ell$, the nonreal eigenvalues have modulus at most $\sqrt2$.
Therefore
\[
 \operatorname{tr}B_J^\ell
 \ge m(2-\theta)^\ell-3|V(J)|2^{\ell/2}.
\]
Apply this to the clean blocks in $J_0$, with $\theta=1/(12x)$.
Their zero deficit, boundary at most $2L$, and \eqref{eq:two-scale-block-volume} verify the hypothesis.
Sum the normalized trace over even lengths $6x\le\ell\le24x$.
The first term contributes at least a fixed positive constant times $m$.
The nonreal error is $o(1)$, since $|V(J_0)|=O(r)$.
The repeated-walk bound is at most
$\tfrac32L^2 2^{-s_0}\mu_0=o(\mu_0)$.
The manuscript's trace-to-cycle calculation consequently gives, for large $r$,
\[
 \mu_0\ge m/16-1.
\]
Together with \eqref{eq:two-scale-many-blocks}, this yields
\[
 \mu_0=\Omega\!\left(\frac{r^{1/20}}{x^4}\right).
\]
But \eqref{eq:two-scale-two-consequences} gives
\[
 \mu_0\le\frac{1+o(1)}4\,r^\kappa.
\]
Since $\kappa<1/100$, these bounds contradict each other for large $r$.
This proves Theorem~\ref{thm:two-scale-polynomial-forest}.

\subsection{Transfer to unrestricted cycle counts}

\begin{corollary}[A polynomial unrestricted cycle error]
\label{cor:two-scale-polynomial-cycles}
For the absolute constant $\kappa$ above,
\[
 \psi(r)\le2^{r-1}+(1+o(1))\frac{2^r}{r^\kappa}.
\]
In particular, $\psi(r)-2^{r-1}=O(2^r/r^\kappa)$.
\end{corollary}

\begin{proof}
Normalize a graph by deleting irrelevant trees and suppressing degree-two vertices.
Rank and cycle count are preserved.
The resulting kernel has at most $2r-2$ vertices and a cycle of length $O(\log r)$.
If necessary, subdivide its edges uniformly to work with a simple graph.
The chosen cycle still has length $O(\log r)$.

For this connected cycle witness, the manuscript's marked cubic reduction has connected marked set $P$ and
$|P|=O(\log r)\le2r/\log^2r$.
Theorem~\ref{thm:two-scale-polynomial-forest} therefore bounds the original forest probability by $1/2+r^{-\kappa}$.
The cycle-flip injection for any witness edge gives
\[
 c_e(G)\le2^{r-2}+2^{r-1}r^{-\kappa}.
\]
Deleting that cyclic edge preserves connectedness and reduces rank by one.
The established edge recurrence, with $E_r=\psi(r)-2^{r-1}$, gives
\[
 E_r\le\max\{0,E_{r-1}+2^{r-1}r^{-\kappa}\}
\]
for all sufficiently large $r$.
Iterating and summing the geometric tail yields
\[
 E_r\le O(1)+\sum_{j\le r}2^{j-1}j^{-\kappa}
      =(1+o(1))2^r r^{-\kappa}.
\]
This proves the corollary.
\end{proof}
